\documentclass[10pt,a4paper]{amsart}
\usepackage[margin=19mm]{geometry}
\usepackage{amsmath,amssymb,mathtools}
\usepackage[scr=rsfs,cal=euler]{mathalfa}
\usepackage{xfrac}
\usepackage[unicode,hidelinks,bookmarks=false]{hyperref}
\newcommand{\ie}{i.e.\ }
\newcommand{\cf}{cf.\ }
\newcommand{\resp}{respectively\ }
\newcommand{\Subsection}{\subsection}
\providecommand{\nobreakdash}{\nobreak\hbox{-}}
\setlength{\emergencystretch}{2em}
\multlinegap0pt
\usepackage{amssymb}
\usepackage{mathtools}
\newtheorem{theorem}{Theorem}
\newtheorem{lemma}{Lemma}
\DeclareMathOperator{\sgn}{sgn}
\title{Real-rootedness of Kazhdan--Lusztig and $Z$-polynomials of thagomizer matroids and graphic matroids of $K_{2,n}$}
\author{Philip B. Zhang}
\date{Source: arXiv:2608.02303v1 (CC BY 4.0)}
\begin{document}
\maketitle

\noindent\emph{Source and licence.} Real-rootedness of Kazhdan--Lusztig and $Z$-polynomials of thagomizer matroids and graphic matroids of $K_{2,n}$. Version arXiv:2608.02303v1 (CC BY 4.0), distributed by arXiv under the Creative Commons Attribution 4.0 International licence, http://creativecommons.org/licenses/by/4.0/, as recorded in the arXiv licence metadata for this exact version and shown on the abstract page https://arxiv.org/abs/2608.02303v1. Full licence notice: \texttt{licenses/arxiv-2608.02303-CC-BY-4.0.txt}.\par\smallskip

\noindent\textit{Editorial scope.} Counted unit: the article's theorem thm:KLpencil (source0.tex lines 99--103). The article's complete original proof of it is delivered with it (source0.tex lines 335--494, one \texttt{proof} environment of 160 lines, which the article places away from the statement). No mathematics is written, repaired or re-derived: the statement and the proof are the article's own lines, in the article's own notation, and no retained line is substituted or re-flowed.  Retained uncounted as the source-local context the proof needs: lines 90--92; lines 193--196; lines 229--231; lines 236--238; lines 248--255.  Nothing was omitted from any retained span, so the package carries no omission record.  No retained line contains a citation, so the wrapper carries no bibliography and no bibliography entry.  No retained line contains a figure environment, an image-inclusion command or drawing code, so no image is delivered and the package declares no asset dependency.  Editorial formatting only, in addition: an AMS \texttt{amsart} preamble that restates the article's own declarations and macros used by the retained lines, namely the environments \texttt{theorem}, \texttt{lemma} on separate counters, together with the standard packages those lines need.  The retained environments print in this document as Theorem 1, Lemma 1 in order of appearance, whereas the article prints the same items with the numbering of its own full document.  The complete native source of the article as distributed in the arXiv e-print for this exact version is bundled unchanged as \texttt{sources/206681/source0.tex}.
\begin{equation}\label{eq:ThagGF}
   \sum_{n\ge0}P_n(x)z^n=C\bigl(z-(1-x)z^2\bigr).
\end{equation}

\begin{equation}\label{eq:PhiUnitCircle}
\Phi_N(A)(e^{2i\theta})
=2^Ne^{iN\theta}\cos^N\theta\,A(-\tan^2\theta).
\end{equation}

\begin{equation}\label{eq:Bgenfun}
\sum_{k\ge0}B_k(q)z^k=C((1+q)z-4qz^2).
\end{equation}

\begin{equation}\label{eq:betaMixedCoeff}
\beta_{k,r}=[X^{k-r}Y^r]\mathcal C(X,Y).
\end{equation}

\begin{lemma}\label{lem:twoVariableCatalan}
One has
\begin{equation}\label{eq:CXYdecomp}
C(X+Y-4XY)
=C(X)+C(Y)-1-\frac{UV(U+V)}{1+UV}.
\end{equation}
Consequently, $[X^aY^b]C(X+Y-4XY)\le0$ whenever $a,b\ge1$.
\end{lemma}

\begin{theorem}\label{thm:KLpencil}
For $n\ge2$ and $0\le\lambda\le n/2$, the polynomial
$P_n(x)+\lambda x$ has exactly $\lfloor n/2\rfloor$ zeros, all of
which are negative and simple.
\end{theorem}

\begin{proof}[Proof of Theorem~\ref{thm:KLpencil}]
Fix $n\ge2$ and $0\le\lambda\le n/2$.  Put
$P_{n,\lambda}(x)=P_n(x)+\lambda x$ and
$d=\lfloor n/2\rfloor$.

\medskip
\noindent\emph{Step 1: Coefficient signs and a root-of-unity bound.}
By \eqref{eq:betaMixedCoeff}, the symmetry of $\mathcal C(X,Y)$ gives
$\beta_{n,r}=\beta_{n,n-r}$.  Moreover,
$\mathcal C(X,0)=C(X)$ and $\mathcal C(0,Y)=C(Y)$, while
Lemma~\ref{lem:twoVariableCatalan} gives
$\beta_{n,0}=\beta_{n,n}=C_n$ and
$\beta_{n,r}\le0$ for $1\le r\le n-1$.
Setting $q=1$ in \eqref{eq:Bgenfun} yields
$\sum_{k\ge0}B_k(1)z^k=C(2z-4z^2)=(1-2z)^{-1}$.  Indeed,
$1-4(2z-4z^2)=(1-4z)^2$, whose formal square root with constant
term $1$ is $1-4z$.  Hence
$B_n(1)=2^n$.

For $j\in\mathbb Z$, let $q=e^{2\pi ij/n}$.  Since $B_n$ has real
coefficients and is palindromic with respect to $n$, we have
$B_n(q)=q^nB_n(q^{-1})=B_n(\overline q)=\overline{B_n(q)}$; hence
$B_n(q)$ is real.  Taking real parts in
$B_n(q)-B_n(1)=\sum_{r=1}^{n-1}\beta_{n,r}(q^r-1)$ gives
\begin{align}
B_n(q)-B_n(1)
&=\sum_{r=1}^{n-1}\beta_{n,r}
\left(\cos\frac{2\pi jr}{n}-1\right) \notag\\
&=2\sum_{r=1}^{n-1}(-\beta_{n,r})
\sin^2\frac{\pi jr}{n}\ge0.
\label{eq:BnRootUnityBound}
\end{align}
Thus $B_n(q)\ge2^n$ for every $n$th root of unity $q$.  This
uniform lower bound is the input needed to control the additional
term arising from $\lambda x$.

\medskip
\noindent\emph{Step 2: Alternating signs on the negative real axis.}
Let $0\le j\le d$ when $n$ is odd, and let $0\le j\le d-1$ when
$n$ is even.  Put
$\theta=j\pi/n$ and $q=e^{2i\theta}$.  By linearity of $\Phi_n$ and
\eqref{eq:PhiUnitCircle},
\begin{equation}\label{eq:pencilTransformed}
2^n(-1)^j\cos^n\theta\,
P_{n,\lambda}(-\tan^2\theta)
=B_n(q)-2^n\lambda(-1)^j
\sin^2\theta\cos^{n-2}\theta.
\end{equation}
If $j$ is odd, the second term on the right-hand side of
\eqref{eq:pencilTransformed} is nonnegative, and the whole expression
is positive by \eqref{eq:BnRootUnityBound}.  If $j$ is even, the
right-hand side of \eqref{eq:pencilTransformed} is at least
$2^n(1-\lambda\sin^2\theta\cos^{n-2}\theta)$.  We claim that
\begin{equation}\label{eq:perturbBound}
0\le\sin^2\theta\cos^{n-2}\theta<\frac2n.
\end{equation}
If $n=2$, then $d=1$ and the only relevant index is $j=0$, so the
left-hand side is $0$.  For $n>2$, put
$u=\cos^2\theta$.  The maximum of
$(1-u)u^{(n-2)/2}$ on $0\le u\le1$ is attained at
$u=(n-2)/n$ and equals
$\frac2n((n-2)/n)^{(n-2)/2}<2/n$.  By
\eqref{eq:perturbBound} and $0\le\lambda\le n/2$, the right-hand
side of \eqref{eq:pencilTransformed} is positive.  Therefore
\begin{equation}\label{eq:pencilSign}
\sgn P_{n,\lambda}\!\left(-\tan^2\frac{j\pi}{n}\right)=(-1)^j.
\end{equation}
These evaluation points lie on the nonpositive real axis and are
strictly decreasing as $j$ increases.  We now lift these signs to the
standard cosine nodes associated with a Chebyshev polynomial; this
supplies the zero-counting step.

\medskip
\noindent\emph{Step 3: Degree, Chebyshev interlacing, and zero counting.}
Expanding $C(t)=\sum_{r\ge0}C_rt^r$ in \eqref{eq:ThagGF} and
extracting the coefficient of $z^n$ gives
\begin{equation}\label{eq:PnExpansion}
P_n(x)=\sum_{m=0}^{\lfloor n/2\rfloor}
C_{n-m}\binom{n-m}{m}(x-1)^m.
\end{equation}
By \eqref{eq:PnExpansion}, if $d\ge2$, then
$[x^d]P_{n,\lambda}(x)=C_{n-d}\binom{n-d}{d}>0$.  If $d=1$, then
$[x]P_{n,\lambda}(x)=(n-1)C_{n-1}+\lambda>0$.  Thus
$\deg P_{n,\lambda}=d$ and its leading coefficient is positive.

Write $P_{n,\lambda}(x)=\sum_{m=0}^d a_mx^m$ and define its
degree-$n$ normalization by
\begin{equation}\label{eq:normalizedP}
\widetilde P_{n,\lambda}(y)
=\sum_{m=0}^d a_m y^{n-2m}(y^2-1)^m.
\end{equation}
For $y\ne0$,
\begin{equation}\label{eq:normalizedPRational}
\widetilde P_{n,\lambda}(y)
=y^nP_{n,\lambda}(1-y^{-2}).
\end{equation}
The polynomial $\widetilde P_{n,\lambda}$ contains only powers
congruent to $n$ modulo $2$, and
\[
[y^n]\widetilde P_{n,\lambda}(y)
=P_{n,\lambda}(1)=C_n+\lambda>0,
\]
where $P_n(1)=C_n$ follows from \eqref{eq:ThagGF}.  Thus its degree
is exactly $n$.

Let $U_m$ denote the Chebyshev polynomial of the second kind,
normalized by
\begin{equation}\label{eq:ChebyshevSecondKind}
U_m(\cos\theta)=\frac{\sin((m+1)\theta)}{\sin\theta},
\end{equation}
and put $c_j=\cos(j\pi/n)$ for $0\le j\le n$.  If
$0\le j<n/2$, then $c_j>0$, and
\eqref{eq:normalizedPRational} together with \eqref{eq:pencilSign}
gives
\[
\sgn\widetilde P_{n,\lambda}(c_j)
=\sgn P_{n,\lambda}\!\left(-\tan^2\frac{j\pi}{n}\right)
=(-1)^j.
\]
The parity of $\widetilde P_{n,\lambda}$ gives the same formula at
the negative nodes.  If $n=2d$ and $j=d$, then
\[
\widetilde P_{n,\lambda}(0)=(-1)^d a_d,
\]
whose sign is $(-1)^d$ by the positivity of the leading coefficient
proved above.  Consequently,
\begin{equation}\label{eq:normalizedNodeSign}
\sgn\widetilde P_{n,\lambda}(c_j)=(-1)^j
\qquad(0\le j\le n).
\end{equation}

Since $c_0>c_1>\cdots>c_n$, the intermediate value theorem gives at
least one zero of $\widetilde P_{n,\lambda}$ in each interval
$(c_j,c_{j-1})$, $1\le j\le n$.  Its degree is $n$, so these are
exactly its zeros and all are simple.  If they are denoted by
$\rho_1>\cdots>\rho_n$, then
\begin{equation}\label{eq:ChebyshevInterlacing}
1>\rho_1>c_1>\rho_2>c_2>\cdots>
c_{n-1}>\rho_n>-1.
\end{equation}
The zeros of $U_{n-1}$ are $c_1,\ldots,c_{n-1}$, so
\eqref{eq:ChebyshevInterlacing} says precisely that $U_{n-1}$
strictly interlaces $\widetilde P_{n,\lambda}$.

It remains to translate this back to $P_{n,\lambda}$.  If $n=2d$,
the $n$ zeros of $\widetilde P_{n,\lambda}$ occur in $d$ pairs
$\{r,-r\}$ with $0<r<1$.  If $n=2d+1$, then $0$ is one of its
zeros and the remaining zeros occur in $d$ such pairs.  By
\eqref{eq:normalizedPRational}, each pair corresponds to a single zero
\[
x=1-r^{-2}<0
\]
of $P_{n,\lambda}$.  Distinct positive values of $r$ give distinct
values of $x$.  Moreover, the change of variables
$x=1-y^{-2}$ has nonzero derivative at $y=r$, so simplicity of the
zeros of $\widetilde P_{n,\lambda}$ implies simplicity of these
zeros of $P_{n,\lambda}$.  We have therefore obtained $d$ distinct
simple negative zeros.  Since $\deg P_{n,\lambda}=d$, these are all
its zeros.
\end{proof}

\end{document}
